\documentclass[11pt,a4paper]{article}

\usepackage[utf8]{inputenc}
\usepackage[T1]{fontenc}
\usepackage{amsmath, amssymb, amsthm}
\usepackage{geometry}
\usepackage{hyperref}
\usepackage{listings}
\usepackage{xcolor}
\usepackage{bm}

\geometry{margin=1in}

\definecolor{codegreen}{rgb}{0,0.6,0}
\definecolor{codegray}{rgb}{0.5,0.5,0.5}
\definecolor{codepurple}{rgb}{0.58,0,0.82}
\definecolor{backcolour}{rgb}{0.95,0.95,0.92}

\lstdefinestyle{mystyle}{
    backgroundcolor=\color{backcolour},   
    commentstyle=\color{codegreen},
    keywordstyle=\color{magenta},
    numberstyle=\tiny\color{codegray},
    stringstyle=\color{codepurple},
    basicstyle=\ttfamily\footnotesize,
    breakatwhitespace=false,         
    breaklines=true,                 
    captionpos=b,                    
    keepspaces=true,                 
    numbers=left,                    
    numbersep=5pt,                  
    showspaces=false,                
    showstringspaces=false,
    showtabs=false,                  
    tabsize=2
}

\lstset{style=mystyle}

\title{Implementation Breakdown: Recursive Asymptotic Variance Estimation for DTMCs}
\author{MS Thesis Technical Documentation}
\date{September 2026}

\begin{document}

\maketitle

\section{Introduction}
This document provides a technical explanation of the Python implementation (\texttt{src/simulate\_dtmc.py}) of the recursive stochastic approximation (SA) algorithm for estimating the asymptotic variance $\kappa(f)$ of a function $f$ defined on a Discrete-Time Markov Chain (DTMC). 

The implementation is based on the research paper: \textit{"Markov Chain Variance Estimation: A Stochastic Approximation Approach"} (Agrawal et al., 2024).

\section{Mathematical Foundations}
Given an irreducible and aperiodic DTMC on a finite state space $\mathcal{S}$ with transition matrix $P$ and stationary distribution $\pi$, we define the stationary expectation of a bounded function $f: \mathcal{S} \to \mathbb{R}$ as $\bar{f} = \mathbb{E}_\pi[f(X)]$.

\subsection{Asymptotic Variance $\kappa(f)$}
The asymptotic variance is defined via the Central Limit Theorem for Markov Chains:
\begin{equation}
    \frac{\sum_{k=0}^{n-1} f(X_k) - n\bar{f}}{\sqrt{n}} \implies \mathcal{N}(0, \kappa(f)) \quad \text{as } n \to \infty
\end{equation}
The implementation uses a formulation involving the solution $V$ to the \textbf{Poisson Equation}:
\begin{equation}
    f - \bar{f}\mathbf{1} = V - PV
\end{equation}
Rearranging the terms, the script evaluates $\kappa(f)$ using the identity:
\begin{equation}
    \kappa(f) = \mathbb{E}_\pi[2f(X)V(X) - 2f(X)\bar{V} - f^2(X) + f(X)\bar{f}]
\end{equation}
where $\bar{V} = \mathbb{E}_\pi[V(X)]$.

\section{Algorithm Implementation}
The core of the simulation is the \texttt{TabularVarianceEstimator} class, which implements \textbf{Algorithm 1} from the paper.

\subsection{The Update Logic}
At each transition $(X_k, X_{k+1})$, the estimator performs four recursive $O(1)$ updates using a step-size $\alpha_k$:
\begin{enumerate}
    \item \textbf{Stationary Mean ($\bar{f}_k$)}: Updates the estimate of $\mathbb{E}_\pi[f]$.
    \item \textbf{Value Function ($V_k$)}: Solves the Poisson equation. The code uses a compact vector update that preserves the orthogonality constraint $\mathbf{1}^T V_k = 0$.
    \item \textbf{Stationary Mean of $V$ ($\bar{V}_k$)}: Tracks the average of the value function.
    \item \textbf{Asymptotic Variance ($\kappa_k$)}: Aggregates the terms from the $\kappa(f)$ identity using linear SA.
\end{enumerate}

\section{Verification and Invariants}
To ensure the simulation is correct, the implementation includes several verification mechanisms:

\subsection{Analytical Ground Truth}
The script includes a \texttt{compute\_analytical\_ground\_truth} function that solves the system exactly using matrix inversion and least squares. It cross-validates the result using two independent formulations of $\kappa(f)$:
\begin{enumerate}
    \item $\kappa(f) = 2\mathbb{E}_\pi[(f(X) - \bar{f})V(X)] - \mathbb{E}_\pi[(f(X) - \bar{f})^2]$
    \item $\kappa(f) = \mathbb{E}_\pi[V^2(X)] - \mathbb{E}_\pi[(PV)^2(X)]$
\end{enumerate}

\subsection{Spectral Gap and Convergence Constants}
The convergence of the algorithm depends on the spectral gap $\Delta_1$ of the symmetric part of $D_\pi(I-P)$. The function \texttt{compute\_theoretical\_constants} uses this gap to derive scale factors $c_1, c_2, c_3$ that satisfy the stability conditions derived in Theorem 2.1 of the paper.

\subsection{Unit Testing}
The \texttt{src/test\_simulate\_dtmc.py} suite verifies the \textbf{orthogonal conservation invariant}. It mathematically proves that for any update $V_{k+1}$, the property $\sum_{i \in \mathcal{S}} V_{k+1}(i) = 0$ is preserved exactly (within machine precision) without requiring an explicit projection step at every iteration.

\section{Usage and Visualization}
The simulation generates a convergence plot (\texttt{src/convergence\_plot.png}) comparing diminishing step-sizes (guaranteed $O(1/n)$ convergence) against constant step-sizes (faster initial descent but higher asymptotic bias).

\end{document}
